\chapter{Configuration and Secrets in the Cluster}
\label{ch:secrets}

\section{Two injection mechanisms}

Part~I established that env vars outrank every file. Kubernetes offers
two ways to set them, and the manifests use both deliberately:

\begin{lstlisting}[language=yaml]
env:                                     # literal, visible in git
  - { name: SPRING_DATASOURCE_URL,
      value: "jdbc:postgresql://postgres-course:5432/ts_course" }
envFrom:                                 # every key of a Secret
  - secretRef: { name: postgres-course-app }
  - secretRef: { name: jwt }
\end{lstlisting}

The split \emph{is} the security boundary: topology (hostnames, ports,
feature flags) is not secret and belongs in reviewable git; credentials
never appear in git at all. When you read a Deployment here, every
\code{env:} entry is public configuration and every \code{envFrom:} line
is a pointer to something created out-of-band.

\section{Secrets: what they are and are not}

A Secret is a namespaced object holding base64-encoded key/value pairs.
Two honest statements about it:

\begin{itemize}
  \item Base64 is \emph{encoding}. \code{echo cm9vdA== | base64 -d}
  reveals the password. The protections are access control (RBAC decides
  who may read Secrets --- notice Chapter~22 grants Jenkins no such
  verb), separate lifecycle from code, and optional encryption at rest.
  \item It is still the right container: every mechanism above it
  (Sealed Secrets, Vault, cloud secret managers) ultimately materializes
  a Secret for pods to consume. Learning the plain object first is
  learning the interface.
\end{itemize}

\section{\code{create-secrets.sh}, annotated}

Secrets here are created by an idempotent script, not manifests:

\begin{lstlisting}[language=bash]
kubectl -n "$NS" create secret generic postgres-course \
  --from-literal=POSTGRES_USER="$PG_USER" \
  --from-literal=POSTGRES_PASSWORD="$PG_PASSWORD" \
  --from-literal=POSTGRES_DB=ts_course \
  --dry-run=client -o yaml | kubectl apply -f -     (*@\co{1}@*)

kubectl -n "$NS" create secret generic postgres-course-app \
  --from-literal=SPRING_DATASOURCE_USERNAME="$PG_USER" \
  --from-literal=SPRING_DATASOURCE_PASSWORD="$PG_PASSWORD" \
  --dry-run=client -o yaml | kubectl apply -f -     (*@\co{2}@*)

kubectl -n "$NS" create secret generic jwt \
  --from-literal=JWT_SECRET="$JWT_SECRET"           (*@\co{3}@*)
\end{lstlisting}

\begin{description}
  \item[\co{1}] The idempotent-create idiom: \code{create} alone fails
  if the Secret exists; rendering with \code{--dry-run} and piping to
  \code{apply} makes the script safely re-runnable --- create or update,
  either way converging. A tiny declarative-over-imperative lesson
  in shell form.
  \item[\co{2}] The same credential twice, spelled per consumer:
  the Postgres image wants \code{POSTGRES\_USER}; Spring wants
  \code{SPRING\_DATASOURCE\_USERNAME}. Two Secrets mean each pod's
  \code{envFrom} imports only names it understands --- no pod carries a
  grab-bag of foreign variables.
  \item[\co{3}] One Secret consumed by two Deployments --- auth-service
  signs JWTs with it, api-gateway validates. The shared-key contract of
  Chapter~5, expressed as infrastructure: rotating it is one script run
  and two pod restarts, and the two consumers \emph{cannot} drift.
\end{description}

\begin{decision}[script, not manifests --- and not in the kustomization]
Secret manifests in git would put base64'd credentials one
\code{base64 -d} away from anyone with repo access, forever, in
history. The script keeps values in the operator's environment
(\code{JWT\_SECRET=... ./create-secrets.sh}), defaults matching dev so
behavior is identical, and git holding only the \emph{shape}. The
kustomization's comment says the same thing from the other side:
secrets are deliberately absent from \code{resources:}.
\end{decision}

\section{ConfigMaps}

The Secret's non-secret twin --- same mechanics, no pretense of
sensitivity, plus a talent Secrets share but ConfigMaps use more:
mounting as \emph{files}. This project has none yet for an honest
reason: its non-secret config already flows through the config-server
and env vars. The natural first ConfigMap arrives with Chapter~26's
config-server retirement --- the \code{configurations/*.yml} files
mounted into pods, letting Spring read them natively and deleting a
service.

\begin{pitfall}[changed the Secret, nothing happened]
Symptom: you re-run \code{create-secrets.sh} with a new password;
services keep using the old one, until some pod restarts days later and
starts failing authentication --- an outage on a delay timer.
Cause: env vars are injected \emph{at container start}; running pods
never see Secret updates. Rotation is therefore two steps by design:
update the Secret, then \code{kubectl -n ts rollout restart deployment
...} for every consumer, immediately --- while old and new credentials
both work if the backing system allows it.
\end{pitfall}

\section{Outside the project}

The gap between this setup and industrial practice is one concept:
making secrets \emph{deliverable by pipeline} without being readable in
git. Three rungs on that ladder: \textbf{Sealed Secrets} (encrypt with
the cluster's public key; commit the sealed blob; a controller decrypts
in-cluster --- GitOps-compatible with one moving part);
\textbf{External Secrets Operator} (a Secret manifest that is a
\emph{pointer} into Vault/AWS Secrets Manager, synced by a controller);
and \textbf{Vault injection} (secrets never become Kubernetes objects at
all). Each replaces the script; none changes how pods consume ---
\code{envFrom} survives every migration, which is why the manifests are
already future-proof.

\begin{quiz}
\begin{enumerate}
  \item Base64 is not encryption --- so name the three things that
  actually protect a Secret, and where this project stands on each.
  \item Why two Secrets for one Postgres credential? What failure does
  the split prevent?
  \item Describe complete JWT-secret rotation on this cluster: commands,
  order, and the failure mode of doing only half.
  \item Which config item in this project would move to a ConfigMap
  first, and what service does that retire?
\end{enumerate}
\end{quiz}
